\section{Implement kMeans with Nested Forward \& Reverse AD}
My code validates and is provided below:
\begin{lstlisting}[language={futhark}]
-- Task 4 (a)
def cost [n] [k] [d] (points: [n][d]f32) (centres: [k][d]f32) : f32 =
  let min_distances =
    map (\p ->
           map (\c -> euclid_dist_2 p c) centres
           |> reduce f32.min f32.inf)
        points
  in reduce (+) 0f32 min_distances

def tolerance = 1 : f32

-- Task 4 (b)
let (cost', cost'') =
  let f centers' = vjp (cost points) centers' 1f32
  in jvp2 f cluster_centers (replicate (k * d) 1f32 |> unflatten)

-- Task 4 (c)
let new_centers =
  let f i j =
    let tmp = cost''[i, j]
    in if tmp == 0f32
       then cluster_centers[i, j]
       else cluster_centers[i
            ,j]
            - (cost'[i, j] / cost''[i, j])
  in map (\i -> map (f i) (iota d)) (iota k)
\end{lstlisting}

\paragraph{Task a:} The cost function simply nests a map inside a map finding the distance from
each point to each center, and then finding the minimum. Once this is done, we
simply reduce over all of the minimum distances to find the cost.The
euclid\_dist\_2 function was provided to us.

\paragraph{Task b:} Here I simply followed the comments explanation. I.e. I
apply the vjp function to the (cost points) in the point that corresponds to
the current cluster center, which is passed to the function f and simply
adjoint with 1. Jvp2 is then applied to this function, which receives the
cluster centers and applies them together with the matrix of ones.

\paragraph{Task c:} This task again follows the comment, but builds upon this
with checking for division by 0. In such a case, we do not update the value of
the center, and in all other cases, we simply provide what the comment states.
This is then mapped to all i and j combinations. 




